\section{Modo JVM ou executável nativo}
\label{sec:jvm-vs-nativo}

O GraalVM Native Image compila a aplicação inteira, antes da execução,
num executável. O guia \citetitle{doc:graalvm-native-basics} dá a
restrição central: depois de construído o binário, nenhum elemento novo
pode ser adicionado em runtime, por exemplo por class loading. A
documentação chama isso de \termo{closed-world assumption}{Suposição de
mundo fechado: tudo o que pode executar precisa ser conhecido no build.
Só os elementos alcançáveis pela análise estática entram no binário.}.
Compare com a seção~\ref{sec:class-loading}: na HotSpot, uma classe pode
ser carregada a qualquer momento; no nativo, não existe class loading
dinâmico, e o \emph{image heap} já traz objetos criados durante o build.

O que muda em relação à Parte~\ref{parte:java-jvm}:
\begin{itemize}
\item O código é compilado antes da execução, não pelo JIT da HotSpot
  em runtime (seção~\ref{sec:jit}). O guia do Quarkus atribui a isso a
  diferença de throughput: em cargas longas, o JIT pode superar o
  nativo.
\item O GC padrão é outro (seção~\ref{sec:memoria-container}).
\item Reflexão, recursos e proxies precisam ser declarados no build
  (seção~\ref{sec:reflexao-nativo}).
\end{itemize}

\subsection*{O que o fornecedor mediu}

O guia \citetitle{doc:quarkus-building-native} publica a tabela abaixo,
atribuída ao benchmark \emph{tuned} do perf-lab do Quarkus de
21/04/2026: Quarkus~3.34.3, JDK~25.0.2, GraalVM~25.0.2-graalce, 4~CPUs,
\texttt{-Xmx512m}. A linha nativa leva o rótulo \enquote{Mandrel},
embora a configuração cite o GraalVM CE; a página não explica a
diferença. O guia não nomeia a aplicação. Os resultados dessa data
ficam em
\url{https://github.com/quarkusio/benchmarks/tree/main/results/spring-quarkus-perf-comparison/2026-04-21_09-26-13},
o que aponta para o repositório \citetitle{doc:quarkus-perf-comparison}:
um serviço REST completo, com JPA/Hibernate sobre PostgreSQL e, na
estratégia \emph{tuned}, métricas e OpenTelemetry. No comando de exemplo
do repositório, a aplicação fica em quatro núcleos fixos
(\texttt{--cpus-app 0,1,2,3}), separada do banco e do gerador de carga.
\textbf{Números do fornecedor.}

\small
\begin{longtable}{@{}p{4cm} p{2.4cm} p{3.6cm} p{1.8cm}@{}}
\toprule
Modo & 1ª requisição & Throughput de pico & RSS \\
\midrule \endhead
JVM (\emph{fast-jar}) & 4417~ms & 13\,265~tps (referência) & 304~MiB \\
JVM + cache AOT (Leyden) & 1859~ms & 12\,389~tps ($\approx$7\% abaixo) &
  240~MiB \\
Nativo (Mandrel) & 581~ms & 5\,411~tps ($\approx$59\% abaixo) & 95~MiB \\
\bottomrule
\end{longtable}
\normalsize

A \enquote{1ª requisição} é o \emph{time to first request}: do comando
de start até a primeira resposta com sucesso, não o \enquote{started
in} do log (guia \citetitle{doc:quarkus-performance-measure}). Isso
explica por que 4,4~s não contradiz o \enquote{sub-second startup} que o
guia \citetitle{doc:quarkus-building-native} atribui ao modo JVM: nessa
página, o \emph{cold start} do
modo JVM vai de cerca de 0,4~s (um REST pequeno) a cerca de 3~s (um CRUD
grande), vindo de outros posts do blog, e a medição do perf-lab inclui
subir o serviço completo e atender a primeira requisição.

Leia a tabela com a Parte~\ref{parte:java-jvm} na cabeça. O nativo ganha em tempo até a
primeira requisição e em memória; o JVM ganha em throughput de pico, que
é onde o JIT, compilando os métodos quentes com o perfil da execução
(seção~\ref{sec:jit}), faz diferença. A linha do meio usa o cache AOT
do Projeto Leyden, a mesma família de recursos dos JEPs~483 e 514
(seção~\ref{sec:class-loading}); o guia não detalha as opções usadas.
É parte do ganho de startup sem sair da HotSpot. Na mesma página, sem atribuir a fonte, o
guia estima o build nativo em 3 a 10 minutos e 4 a 8~GB de RAM na
máquina de build. Tamanhos de imagem e tempos de \emph{cold start} da
página vêm de outros posts do blog, com outras condições; por isso não
entram na tabela.

\subsection*{A decisão}

O guia resume: \enquote{Start with JVM mode. Move to native when you
have a concrete need.}

\begin{description}
\item[Quando escolher JVM] Serviço de longa duração em que o throughput
  de pico importa; build rápido no ciclo de desenvolvimento e no CI; uso
  intenso de reflexão ou class loading dinâmico, inclusive por
  bibliotecas sem extensão Quarkus.
\item[Quando escolher nativo] Startup é crítico (serverless, CLI,
  \emph{scale-to-zero}, segundo o guia); memória por réplica é o custo
  dominante; você quer a menor imagem possível.
\item[O que medir] Seguindo o guia \citetitle{doc:quarkus-performance-measure}:
  \emph{time to first request} (do comando de start até a primeira
  resposta 200, não o tempo que o framework imprime no log), RSS com
  \texttt{ps} ou \texttt{pmap} (o guia pede explicitamente para
  \emph{não} usar \texttt{docker stats}), throughput e p99 sob o teste
  de carga da Parte~\ref{parte:testes}, e o tempo de build no CI. Compare as três
  opções da tabela no \emph{seu} serviço antes de decidir.
\end{description}

\paragraph{Mandrel ou GraalVM.}
O mesmo guia descreve o Mandrel como uma distribuição \emph{downstream}
do GraalVM CE, construída sobre o OpenJDK padrão e sem suporte a
programação poliglota, e o recomenda para alvos Linux em container.
Para builds com \texttt{-Dnative} sem GraalVM instalado, há
\texttt{-Dquarkus.native.container-build=true}.
